\documentclass[11pt,a4paper]{article}

\usepackage[a4paper,margin=2.5cm]{geometry}
\setlength{\parindent}{0pt}
\setlength{\parskip}{0.6em}
\usepackage{helvet}
\renewcommand{\familydefault}{\sfdefault}

\usepackage[utf8]{inputenc}
\usepackage[T1]{fontenc}
\usepackage{setspace}
\onehalfspacing
\usepackage{graphicx}

\usepackage{titlesec}
\titleformat{\section}{\bfseries\Large\centering}{}{0pt}{}
\titlespacing*{\section}{0pt}{2em}{1em}

\titleformat{\subsection}{\bfseries\large}{}{0pt}{}
\titlespacing*{\subsection}{0pt}{1em}{0.5em}

\begin{document}

\begin{center}
\includegraphics[angle=90,width=3cm]{logo_preta.png}\\[1em]
\textbf{\Large ESTATUTO SOCIAL DO MOVIMENTO EM DEFESA DA EDUCAÇÃO}\\
\textbf{\large MoveEduca}
\end{center}

\section*{Nota Técnica}

O presente estatuto foi reestruturado para reduzir riscos de exigências cartorárias, nulidades e conflitos internos, com correção da numeração, eliminação de duplicidades, definição de quóruns, órgãos, competências, regras de eleição, vacância, responsabilidade, conflito de interesses, prestação de contas e destinação patrimonial. Foram incluídas cláusulas compatíveis com o Código Civil, a Lei nº 13.019/2014, a Lei de Registros Públicos, as Normas Brasileiras de Contabilidade aplicáveis às entidades sem fins lucrativos e as exigências usuais para atuação como Organização da Sociedade Civil em parcerias com a Administração Pública.

\newpage

\section*{TÍTULO I -- DA DENOMINAÇÃO, NATUREZA JURÍDICA, SEDE, FORO E DURAÇÃO}

\textbf{Art. 1º} O \textbf{Movimento em Defesa da Educação}, com denominação abreviada \textbf{MoveEduca}, é uma associação civil, pessoa jurídica de direito privado, sem fins lucrativos, de fins não econômicos, apartidária, com finalidade educacional e social, constituída em \textbf{[DATA DA ASSEMBLEIA]}, nos termos dos arts. 44, 46 e 53 a 61 do Código Civil.

\textbf{§ 1º} A associação terá sede e foro na cidade de \textbf{[CIDADE/UF]}, no endereço \textbf{[ENDEREÇO COMPLETO]}, CEP \textbf{[CEP]}.

\textbf{§ 2º} A associação terá duração por prazo indeterminado.

\textbf{§ 3º} Mediante deliberação da Diretoria Executiva, observadas as diretrizes do Conselho de Administração e as normas legais aplicáveis, a associação poderá criar, manter, transferir ou extinguir filiais, escritórios, núcleos, representações, unidades de projeto ou bases operacionais em qualquer localidade do território nacional ou no exterior.

\textbf{§ 4º} A denominação alternativa \textbf{Move\&Educa} poderá ser utilizada exclusivamente como nome fantasia, marca, identidade visual ou denominação de comunicação institucional, sem substituir a denominação civil oficial da associação.

\section*{TÍTULO II -- DAS FINALIDADES, OBJETIVOS E PRINCÍPIOS}

\textbf{Art. 2º} O MoveEduca tem por finalidade principal defender, promover, apoiar e desenvolver a educação, a cidadania, a inclusão social, a cultura, a pesquisa, a tecnologia, a inovação, a formação profissional, a assistência educacional e o fortalecimento institucional de organizações e iniciativas voltadas ao interesse público.

\textbf{Art. 3º} Para a consecução de suas finalidades de relevância pública e social, o MoveEduca poderá:

\textbf{I} -- promover, desenvolver, realizar, apoiar e divulgar estudos, pesquisas, projetos, programas e ações de educação, ensino, formação, capacitação, cidadania, cultura, ciência, tecnologia e inovação;

\textbf{II} -- defender o direito de acesso à educação de qualidade, inclusiva, plural, democrática e socialmente referenciada;

\textbf{III} -- realizar cursos, oficinas, treinamentos, seminários, congressos, conferências, encontros, fóruns, campanhas, publicações e demais atividades educacionais, culturais, científicas, tecnológicas e sociais;

\textbf{IV} -- desenvolver, executar, administrar, monitorar e avaliar programas e projetos voltados à educação, inclusão digital, inclusão produtiva, formação profissional, desenvolvimento comunitário, promoção da cidadania e fortalecimento de capacidades institucionais;

\textbf{V} -- produzir, editar, publicar, distribuir e divulgar conteúdos por meios impressos, digitais, audiovisuais, científicos, técnicos, didáticos, jornalísticos ou institucionais, inclusive livros, artigos, relatórios, manuais, revistas, vídeos, documentários, softwares, plataformas, materiais formativos e recursos educacionais;

\textbf{VI} -- prestar assessoria, consultoria, apoio técnico, formação, mentoria e serviços especializados compatíveis com suas finalidades, inclusive a instituições de ensino, organizações da sociedade civil, entidades públicas e privadas, redes comunitárias e coletivos sociais;

\textbf{VII} -- fortalecer organizações da sociedade civil por meio de apoio em governança, planejamento, mobilização de recursos, transparência, integridade, prestação de contas, tecnologia, gestão de projetos e processos institucionais;

\textbf{VIII} -- atuar, direta ou indiretamente, em iniciativas voltadas a crianças, adolescentes, jovens, adultos, idosos, mulheres, pessoas com deficiência, comunidades vulnerabilizadas, profissionais da educação, lideranças comunitárias e demais públicos relacionados às suas finalidades;

\textbf{IX} -- promover a ética, a paz, a cidadania, os direitos humanos, a democracia, a participação social, o voluntariado, a diversidade, a equidade e a não discriminação;

\textbf{X} -- apoiar ações de sustentabilidade, responsabilidade socioambiental e desenvolvimento sustentável quando relacionadas ou complementares às suas finalidades educacionais e sociais;

\textbf{XI} -- celebrar termos de colaboração, termos de fomento, acordos de cooperação, convênios, contratos, termos de parceria, ajustes e instrumentos congêneres com pessoas jurídicas de direito público ou privado, nacionais ou estrangeiras, conforme a legislação aplicável;

\textbf{XII} -- receber, administrar e aplicar recursos públicos e privados, doações, subvenções, patrocínios, apoios, emendas, editais, contribuições, bens, direitos e receitas compatíveis com seus objetivos estatutários;

\textbf{XIII} -- atuar judicial ou extrajudicialmente, quando cabível e observada a legislação aplicável, na defesa de direitos difusos, coletivos ou individuais homogêneos relacionados à educação e às finalidades institucionais.

\textbf{Art. 4º} No desenvolvimento de suas atividades, o MoveEduca observará os princípios da legalidade, impessoalidade, moralidade, publicidade, economicidade, eficiência, transparência, participação social, equidade, responsabilidade, sustentabilidade, proteção de dados pessoais, integridade e respeito aos direitos humanos.

\textbf{Art. 5º} O MoveEduca não fará discriminação de raça, cor, etnia, nacionalidade, origem, sexo, gênero, identidade de gênero, orientação sexual, idade, deficiência, condição socioeconômica, convicção religiosa, opinião política ou qualquer outra forma de discriminação ilícita.

\textbf{Art. 6º} O MoveEduca não distribuirá entre associados, instituidores, fundadores, dirigentes, conselheiros, empregados, voluntários, mantenedores, doadores ou terceiros eventuais excedentes operacionais, sobras, resultados, dividendos, bonificações, participações ou parcelas de seu patrimônio, sob qualquer forma ou pretexto.

\textbf{Parágrafo único.} Todos os recursos, receitas, rendas, superávits e resultados do MoveEduca serão integralmente aplicados na consecução de suas finalidades estatutárias, admitida a constituição de reservas técnicas, fundos patrimoniais ou fundos de projeto, desde que compatíveis com a legislação aplicável e aprovados pelo órgão competente.

\section*{TÍTULO III -- DOS ASSOCIADOS}

\textbf{Art. 7º} O MoveEduca será constituído por número ilimitado de associados, pessoas naturais maiores de 18 anos ou pessoas jurídicas, admitidos na forma deste Estatuto, que concordem com suas finalidades, normas estatutárias e regimentais.

\textbf{Art. 8º} O quadro associativo será composto pelas seguintes categorias:

\textbf{I} -- associados fundadores: aqueles que participaram da assembleia de constituição e assinaram a respectiva ata;

\textbf{II} -- associados efetivos: aqueles admitidos após a constituição, com direito de votar e ser votados, observados este Estatuto e o Regimento Interno;

\textbf{III} -- associados contribuintes: aqueles que contribuem financeiramente, tecnicamente, voluntariamente ou por outros meios para as atividades da associação, podendo ter direito de votar e ser votados se assim aprovado no ato de admissão ou no Regimento Interno;

\textbf{IV} -- associados honorários ou beneméritos: pessoas físicas ou jurídicas reconhecidas pela Assembleia Geral em razão de relevantes serviços prestados à educação, à sociedade ou ao MoveEduca, sem obrigação de contribuição financeira e sem direito automático de votar ou ser votado, salvo deliberação expressa da Assembleia Geral.

\textbf{Art. 9º} A admissão de associados dependerá de requerimento escrito do interessado, aceitação expressa deste Estatuto e do Regimento Interno, análise pela Diretoria Executiva e aprovação por maioria simples da Diretoria Executiva, com posterior comunicação à Assembleia Geral.

\textbf{§ 1º} O Regimento Interno poderá estabelecer documentos, formulários, critérios objetivos, prazos e procedimentos complementares para admissão de associados.

\textbf{§ 2º} A recusa de admissão deverá ser fundamentada, cabendo pedido de reconsideração à Diretoria Executiva no prazo de 15 (quinze) dias contados da ciência do interessado.

\textbf{Art. 10.} São direitos dos associados em pleno gozo de seus direitos estatutários, conforme sua categoria:

\textbf{I} -- participar das Assembleias Gerais;

\textbf{II} -- votar e ser votado para cargos eletivos, quando a categoria associativa assim permitir;

\textbf{III} -- apresentar propostas, requerimentos, sugestões e pedidos de informação aos órgãos da associação;

\textbf{IV} -- participar de atividades, programas e projetos da associação, observadas as normas internas e as condições de cada iniciativa;

\textbf{V} -- requerer a convocação de Assembleia Geral, na forma deste Estatuto;

\textbf{VI} -- solicitar seu desligamento voluntário a qualquer tempo.

\textbf{Art. 11.} São deveres dos associados:

\textbf{I} -- cumprir este Estatuto, o Regimento Interno, as políticas institucionais e as deliberações regularmente aprovadas pelos órgãos competentes;

\textbf{II} -- zelar pela reputação, patrimônio, finalidade e integridade institucional do MoveEduca;

\textbf{III} -- agir com boa-fé, lealdade, urbanidade, respeito, cooperação e responsabilidade;

\textbf{IV} -- manter atualizados seus dados cadastrais;

\textbf{V} -- cumprir as contribuições, encargos ou compromissos assumidos perante a associação, quando houver;

\textbf{VI} -- declarar situações de conflito de interesses de que tenham conhecimento;

\textbf{VII} -- preservar a confidencialidade de informações internas, estratégicas, pessoais ou sigilosas, salvo autorização expressa ou obrigação legal de divulgação;

\textbf{VIII} -- comunicar à Diretoria Executiva indícios de irregularidades, desvios, violações estatutárias ou riscos relevantes à associação.

\textbf{Art. 12.} O associado poderá desligar-se voluntariamente mediante comunicação escrita à Diretoria Executiva, sem prejuízo da responsabilidade por obrigações assumidas antes do desligamento.

\textbf{Art. 13.} A exclusão de associado somente poderá ocorrer por justa causa, assegurados procedimento escrito, contraditório, ampla defesa e recurso à Assembleia Geral.

\textbf{§ 1º} Considera-se justa causa, entre outras hipóteses: violação grave deste Estatuto ou do Regimento Interno; prática de ato ilícito ou incompatível com as finalidades da associação; dano material ou reputacional relevante; uso indevido do nome, marca ou recursos da associação; inadimplência injustificada e reiterada de obrigações associativas; conflito de interesses não declarado; fraude; assédio; discriminação; ou conduta que comprometa a integridade institucional.

\textbf{§ 2º} A ausência reiterada e injustificada em Assembleias ou atividades associativas poderá ensejar advertência, suspensão de direitos ou exclusão apenas se houver previsão objetiva em Regimento Interno e sempre mediante procedimento formal.

\textbf{§ 3º} O procedimento de exclusão será instaurado pela Diretoria Executiva, de ofício ou mediante provocação fundamentada, com notificação do associado para apresentar defesa no prazo mínimo de 15 (quinze) dias.

\textbf{§ 4º} A decisão de exclusão caberá à Diretoria Executiva, por maioria absoluta de seus membros, devendo ser fundamentada e comunicada por escrito ao associado.

\textbf{§ 5º} Da decisão de exclusão caberá recurso, com efeito suspensivo, à Assembleia Geral, no prazo de 15 (quinze) dias contados da ciência da decisão.

\textbf{Art. 14.} Os associados não respondem, nem mesmo subsidiariamente, pelas obrigações da associação, salvo quando agirem com dolo, fraude, abuso de direito, desvio de finalidade, confusão patrimonial ou nas demais hipóteses previstas em lei.

\section*{TÍTULO IV -- DA ORGANIZAÇÃO, ADMINISTRAÇÃO E GOVERNANÇA}

\textbf{Art. 15.} São órgãos do MoveEduca:

\textbf{I} -- Assembleia Geral;

\textbf{II} -- Diretoria Executiva;

\textbf{III} -- Conselho Fiscal;

\textbf{IV} -- Conselho de Administração.

\textbf{Art. 16.} Os membros da Diretoria Executiva, do Conselho Fiscal e do Conselho de Administração serão eleitos pela Assembleia Geral, com mandato de 2 (dois) anos, permitida 1 (uma) reeleição consecutiva para o mesmo cargo.

\textbf{§ 1º} Após o exercício de 2 (dois) mandatos consecutivos no mesmo cargo, o dirigente ou conselheiro somente poderá voltar a ocupá-lo após intervalo mínimo de 1 (um) mandato.

\textbf{§ 2º} A posse ocorrerá na própria Assembleia de eleição ou em data por ela fixada.

\textbf{§ 3º} Os mandatos poderão ser prorrogados automaticamente, por até 90 (noventa) dias, quando necessário para conclusão do processo eleitoral, prestação de contas ou registro cartorário, desde que não haja oposição da Assembleia Geral.

\textbf{Art. 17.} Em caso de vacância, renúncia, impedimento definitivo ou destituição de membro eleito:

\textbf{I} -- na Presidência, assumirá o Vice-Presidente até o término do mandato;

\textbf{II} -- nos demais cargos da Diretoria Executiva, a Diretoria Executiva indicará substituto interino, submetendo a indicação à ratificação da primeira Assembleia Geral subsequente;

\textbf{III} -- no Conselho Fiscal ou no Conselho de Administração, assumirá suplente, se houver; inexistindo suplente, a primeira Assembleia Geral subsequente elegerá substituto para completar o mandato.

\textbf{Art. 18.} Os dirigentes e conselheiros não respondem pessoalmente pelas obrigações regularmente contraídas em nome da associação, salvo quando agirem com dolo, fraude, culpa grave, abuso de poder, desvio de finalidade, confusão patrimonial, violação da lei, deste Estatuto ou de deliberação regularmente aprovada.

\textbf{Art. 19.} A governança do MoveEduca observará segregação mínima de funções entre execução, aprovação e fiscalização, de modo que:

\textbf{I} -- a Diretoria Executiva executará a gestão administrativa, financeira e operacional;

\textbf{II} -- o Conselho de Administração exercerá orientação estratégica, supervisão institucional e aprovação de matérias de sua competência;

\textbf{III} -- o Conselho Fiscal fiscalizará a gestão econômico-financeira, contábil e patrimonial;

\textbf{IV} -- a Assembleia Geral exercerá as competências soberanas previstas neste Estatuto e na lei.

\textbf{Art. 20.} A associação poderá realizar Assembleias, reuniões da Diretoria Executiva, reuniões do Conselho Fiscal, reuniões do Conselho de Administração, votações, consultas, registros de presença e deliberações por meio presencial, híbrido ou eletrônico, inclusive por aplicativo, plataforma digital ou sistema próprio.

\textbf{§ 1º} A utilização de meio eletrônico deverá assegurar, no mínimo, a identificação dos participantes, a segurança do voto, o direito de participação, o direito de manifestação, a integridade dos registros e a geração de registros auditáveis.

\textbf{§ 2º} Os registros eletrônicos de presença, votos, manifestações e deliberações poderão compor atas, listas de presença, relatórios ou documentos equivalentes, observadas a legislação aplicável, as exigências do Registro Civil de Pessoas Jurídicas e as políticas internas da associação.

\textbf{§ 3º} O Regimento Interno poderá disciplinar procedimentos complementares para convocação, autenticação, votação, registro de presença, guarda de evidências e uso de sistemas eletrônicos.

\section*{TÍTULO V -- DA ASSEMBLEIA GERAL}

\textbf{Art. 21.} A Assembleia Geral é o órgão soberano da associação, constituída pelos associados com direito a voto e em pleno gozo de seus direitos estatutários.

\textbf{Art. 22.} Compete privativamente à Assembleia Geral:

\textbf{I} -- eleger e empossar os membros da Diretoria Executiva, do Conselho Fiscal e do Conselho de Administração;

\textbf{II} -- destituir administradores e conselheiros;

\textbf{III} -- aprovar as contas anuais, demonstrações financeiras, relatório de atividades e parecer do Conselho Fiscal;

\textbf{IV} -- alterar este Estatuto;

\textbf{V} -- deliberar sobre dissolução, liquidação e destinação do patrimônio remanescente;

\textbf{VI} -- aprovar o Regimento Interno e suas alterações, quando submetidos pela Diretoria Executiva ou pelo Conselho de Administração;

\textbf{VII} -- deliberar sobre alienação, oneração, permuta ou aquisição relevante de bens imóveis, bem como sobre operações patrimoniais extraordinárias;

\textbf{VIII} -- julgar recurso contra exclusão de associado;

\textbf{IX} -- deliberar sobre matérias relevantes submetidas pela Diretoria Executiva, pelo Conselho Fiscal, pelo Conselho de Administração ou por associados legitimados;

\textbf{X} -- autorizar a associação a requerer, manter ou renunciar a qualificações, certificações ou títulos públicos cuja legislação exija deliberação assemblear.

\textbf{Art. 23.} A Assembleia Geral reunir-se-á:

\textbf{I} -- ordinariamente, uma vez por ano, preferencialmente até 30 de abril, para apreciar relatório de atividades, demonstrações financeiras, parecer do Conselho Fiscal, contas da administração e, quando aplicável, eleger órgãos estatutários;

\textbf{II} -- extraordinariamente, sempre que necessário.

\textbf{Art. 24.} A Assembleia Geral poderá ser convocada:

\textbf{I} -- pelo Presidente;

\textbf{II} -- pela Diretoria Executiva;

\textbf{III} -- pelo Conselho Fiscal;

\textbf{IV} -- pelo Conselho de Administração;

\textbf{V} -- por requerimento de, no mínimo, 1/5 (um quinto) dos associados com direito a voto;

\textbf{VI} -- por requerimento de 3 (três) associados com direito a voto, quando esse número for inferior a 1/5 (um quinto) do quadro associativo.

\textbf{Art. 25.} A convocação será feita por edital, correio eletrônico, mensagem eletrônica, publicação no sítio eletrônico da associação, correspondência ou outro meio idôneo previsto em Regimento Interno, com antecedência mínima de 5 (cinco) dias, indicando data, horário, local ou meio eletrônico de participação e ordem do dia.

\textbf{§ 1º} Para alteração estatutária, destituição de administradores ou conselheiros e dissolução da associação, a convocação deverá indicar expressamente a matéria e observar antecedência mínima de 10 (dez) dias.

\textbf{§ 2º} A Assembleia somente poderá deliberar sobre matérias constantes da ordem do dia, salvo deliberações de encaminhamento sem caráter decisório.

\textbf{Art. 26.} A Assembleia Geral instalar-se-á:

\textbf{I} -- em primeira convocação, com a presença da maioria absoluta dos associados com direito a voto;

\textbf{II} -- em segunda convocação, 30 (trinta) minutos após o horário designado, com qualquer número de associados com direito a voto, salvo quórum legal ou estatutário especial.

\textbf{Art. 27.} As deliberações gerais da Assembleia serão tomadas por maioria simples dos votos dos associados presentes, salvo disposição legal ou estatutária diversa.

\textbf{§ 1º} Para alteração deste Estatuto e destituição de administradores ou conselheiros, será exigido voto favorável de 2/3 (dois terços) dos associados presentes à Assembleia especialmente convocada para esse fim.

\textbf{§ 2º} Para dissolução da associação, será exigido voto favorável de 2/3 (dois terços) dos associados com direito a voto presentes à Assembleia especialmente convocada para esse fim, instalada em primeira convocação com maioria absoluta dos associados com direito a voto ou, em segunda convocação, com ao menos 1/3 (um terço) dos associados com direito a voto.

\textbf{§ 3º} Para alienação, oneração, permuta ou aquisição relevante de bens imóveis, será exigido voto favorável da maioria absoluta dos associados presentes, após parecer do Conselho Fiscal e manifestação do Conselho de Administração.

\textbf{Art. 28.} As Assembleias serão registradas em ata, assinada pelo Presidente e pelo Secretário da mesa, podendo ser assinada física ou eletronicamente pelos presentes, conforme a legislação aplicável e as exigências do Registro Civil de Pessoas Jurídicas.

\textbf{Parágrafo único.} Quando a Assembleia ocorrer por meio híbrido ou eletrônico, a ata poderá ser acompanhada de relatório de presença, extrato de votação, gravação, registro de sistema, lista eletrônica ou outro meio idôneo que demonstre participação, manifestação e deliberação.

\section*{TÍTULO VI -- DA DIRETORIA EXECUTIVA}

\textbf{Art. 29.} A Diretoria Executiva é o órgão de administração, representação, gestão e execução das atividades do MoveEduca.

\textbf{Art. 30.} A Diretoria Executiva será composta pelos seguintes cargos:

\textbf{I} -- Presidente;

\textbf{II} -- Vice-Presidente;

\textbf{III} -- Diretor Financeiro;

\textbf{IV} -- Secretário-Geral.

\textbf{Art. 31.} Compete à Diretoria Executiva:

\textbf{I} -- cumprir e fazer cumprir este Estatuto, o Regimento Interno, as políticas institucionais e as deliberações da Assembleia Geral;

\textbf{II} -- administrar a associação e executar seu planejamento, programas, projetos, atividades e parcerias;

\textbf{III} -- elaborar e submeter à Assembleia Geral o relatório anual de atividades, demonstrações financeiras, prestação de contas e proposta de planejamento anual;

\textbf{IV} -- elaborar orçamento, planos de trabalho, propostas, projetos e instrumentos necessários à captação e execução de recursos;

\textbf{V} -- celebrar termos de colaboração, termos de fomento, acordos de cooperação, convênios, contratos, termos de parceria, ajustes e instrumentos congêneres, observadas as competências da Assembleia Geral e do Conselho de Administração;

\textbf{VI} -- admitir, desligar, contratar, supervisionar e remunerar empregados, colaboradores, consultores, prestadores de serviços, voluntários, bolsistas ou profissionais vinculados a projetos, observada a legislação aplicável;

\textbf{VII} -- aprovar despesas, compras, contratações e pagamentos dentro dos limites de alçada definidos neste Estatuto, no Regimento Interno ou em políticas internas;

\textbf{VIII} -- manter escrituração contábil, controles internos, arquivos, documentos, cadastros e registros institucionais;

\textbf{IX} -- instaurar e decidir procedimentos disciplinares de associados, ressalvado recurso à Assembleia Geral;

\textbf{X} -- propor ao Conselho de Administração políticas internas de integridade, compras e contratações, voluntariado, proteção de dados, gestão de riscos, prestação de contas e transparência;

\textbf{XI} -- deliberar sobre criação de filiais, escritórios, núcleos, representações e unidades de projeto;

\textbf{XII} -- praticar os demais atos necessários à gestão ordinária da associação.

\textbf{Art. 32.} A Diretoria Executiva reunir-se-á ordinariamente ao menos uma vez por trimestre e extraordinariamente sempre que necessário, por convocação do Presidente ou da maioria de seus membros.

\textbf{§ 1º} As reuniões da Diretoria Executiva instalar-se-ão com a presença da maioria de seus membros e deliberarão por maioria simples dos presentes.

\textbf{§ 2º} Em caso de empate, prevalecerá o voto do Presidente.

\textbf{Art. 33.} Compete ao Presidente:

\textbf{I} -- representar o MoveEduca ativa e passivamente, judicial e extrajudicialmente;

\textbf{II} -- convocar e presidir reuniões da Diretoria Executiva e Assembleias Gerais, salvo deliberação em contrário da própria Assembleia;

\textbf{III} -- assinar documentos institucionais, contratos, instrumentos de parceria, procurações, correspondências oficiais e atos de representação da associação;

\textbf{IV} -- autorizar despesas, pagamentos e movimentações financeiras, em conjunto com o Diretor Financeiro ou outro dirigente formalmente designado;

\textbf{V} -- constituir procuradores, com poderes específicos e prazo determinado, quando necessário;

\textbf{VI} -- coordenar a execução das deliberações dos órgãos estatutários;

\textbf{VII} -- apresentar relatórios institucionais aos órgãos competentes.

\textbf{Art. 34.} Compete ao Vice-Presidente:

\textbf{I} -- substituir o Presidente em suas faltas, ausências, impedimentos temporários ou vacância, conforme este Estatuto;

\textbf{II} -- apoiar a coordenação institucional, articulação de parcerias, representação externa e acompanhamento de projetos;

\textbf{III} -- exercer atribuições delegadas pelo Presidente ou pela Diretoria Executiva.

\textbf{Art. 35.} Compete ao Diretor Financeiro:

\textbf{I} -- coordenar a gestão financeira, orçamentária, patrimonial e contábil da associação;

\textbf{II} -- arrecadar e controlar contribuições, doações, receitas, auxílios, subvenções, patrocínios, rendimentos e demais entradas financeiras;

\textbf{III} -- efetuar pagamentos, movimentações bancárias e aplicações financeiras autorizadas, em conjunto com o Presidente ou outro dirigente formalmente designado;

\textbf{IV} -- manter numerário, contas bancárias e aplicações financeiras em instituições autorizadas a funcionar no País;

\textbf{V} -- organizar documentos financeiros, fiscais, contábeis, bancários e patrimoniais;

\textbf{VI} -- apresentar relatórios financeiros à Diretoria Executiva, ao Conselho Fiscal, ao Conselho de Administração e à Assembleia Geral, quando solicitado;

\textbf{VII} -- providenciar, com apoio técnico contábil, a elaboração das demonstrações financeiras e prestações de contas.

\textbf{Art. 36.} Compete ao Secretário-Geral:

\textbf{I} -- organizar a secretaria institucional, cadastro de associados, arquivos, livros, atas e documentos da associação;

\textbf{II} -- secretariar Assembleias e reuniões, lavrando ou providenciando as respectivas atas;

\textbf{III} -- expedir editais, comunicações, correspondências, certidões internas e documentos administrativos;

\textbf{IV} -- receber requerimentos de admissão, desligamento e manifestações de associados;

\textbf{V} -- manter atualizados os registros estatutários, regimentais e institucionais;

\textbf{VI} -- substituir, temporariamente, o Vice-Presidente em seus impedimentos, quando designado pela Diretoria Executiva.

\textbf{Parágrafo único.} Na ausência do Secretário-Geral, o Presidente poderá indicar secretário \textit{ad hoc} para secretariar reunião ou Assembleia específica.

\textbf{Art. 37.} A movimentação bancária, emissão de cheques, transferências, contratação de obrigações financeiras e autorização de pagamentos dependerão da assinatura ou autorização conjunta do Presidente e do Diretor Financeiro.

\textbf{§ 1º} Em caso de ausência ou impedimento de um dos signatários, a Diretoria Executiva poderá designar outro dirigente para assinatura conjunta, mediante registro em ata.

\textbf{§ 2º} Despesas ordinárias, recorrentes ou vinculadas a projetos aprovados poderão seguir limites de alçada definidos em Regimento Interno, política financeira, plano de trabalho ou deliberação específica.

\textbf{§ 3º} A aquisição, alienação, oneração, permuta ou promessa de alienação de bens imóveis dependerá de autorização da Assembleia Geral, precedida de parecer do Conselho Fiscal e manifestação do Conselho de Administração.

\section*{TÍTULO VII -- DO CONSELHO DE ADMINISTRAÇÃO}

\textbf{Art. 38.} O Conselho de Administração é órgão colegiado de orientação estratégica, supervisão institucional e apoio à governança, sem função executiva ordinária.

\textbf{Art. 39.} O Conselho de Administração será composto por 3 (três) membros titulares, eleitos pela Assembleia Geral, admitida a eleição de suplentes.

\textbf{§ 1º} Os membros do Conselho de Administração não poderão exercer simultaneamente cargo na Diretoria Executiva.

\textbf{§ 2º} O Conselho de Administração elegerá, entre seus membros, um Coordenador.

\textbf{Art. 40.} Compete ao Conselho de Administração:

\textbf{I} -- acompanhar a execução da missão, das finalidades e das diretrizes estratégicas da associação;

\textbf{II} -- opinar sobre planejamento estratégico, orçamento anual, planos de expansão, políticas institucionais e gestão de riscos;

\textbf{III} -- recomendar boas práticas de governança, integridade, transparência, proteção de dados, compras, contratações, voluntariado e prestação de contas;

\textbf{IV} -- manifestar-se previamente sobre operações patrimoniais relevantes, empréstimos, alienação ou oneração de bens imóveis e parcerias de alto impacto institucional;

\textbf{V} -- acompanhar, sem prejuízo das competências do Conselho Fiscal, a execução de projetos estratégicos e parcerias públicas relevantes;

\textbf{VI} -- propor à Assembleia Geral alterações estatutárias e diretrizes institucionais;

\textbf{VII} -- convocar Assembleia Geral Extraordinária, quando necessário;

\textbf{VIII} -- exercer outras atribuições consultivas ou deliberativas expressamente previstas neste Estatuto, no Regimento Interno ou atribuídas pela Assembleia Geral.

\textbf{Art. 41.} O Conselho de Administração reunir-se-á ordinariamente ao menos 2 (duas) vezes por ano e extraordinariamente sempre que necessário, por convocação de seu Coordenador, da maioria de seus membros, da Diretoria Executiva ou do Conselho Fiscal.

\textbf{Parágrafo único.} As reuniões instalar-se-ão com a presença da maioria de seus membros e deliberarão por maioria simples dos presentes.

\section*{TÍTULO VIII -- DO CONSELHO FISCAL}

\textbf{Art. 42.} O Conselho Fiscal é órgão de fiscalização econômico-financeira, contábil, patrimonial e de prestação de contas da associação.

\textbf{Art. 43.} O Conselho Fiscal será composto por 3 (três) membros titulares, eleitos pela Assembleia Geral, admitida a eleição de suplentes.

\textbf{§ 1º} Os membros do Conselho Fiscal não poderão exercer simultaneamente cargo na Diretoria Executiva ou no Conselho de Administração.

\textbf{§ 2º} O Conselho Fiscal poderá contar com apoio técnico contábil, jurídico, auditoria independente ou especialistas, quando necessário e aprovado pelo órgão competente.

\textbf{Art. 44.} Compete ao Conselho Fiscal:

\textbf{I} -- examinar livros, registros, documentos, extratos, contratos, instrumentos de parceria, prestações de contas e demonstrações financeiras;

\textbf{II} -- emitir parecer sobre balanço patrimonial, demonstrações financeiras, relatório anual de atividades, contas da administração e prestações de contas;

\textbf{III} -- requisitar à Diretoria Executiva documentos e esclarecimentos sobre operações econômico-financeiras, contábeis e patrimoniais;

\textbf{IV} -- acompanhar a aplicação de recursos públicos e privados, especialmente quando vinculados a instrumentos de parceria, editais, planos de trabalho ou projetos específicos;

\textbf{V} -- recomendar medidas de controle interno, correção, transparência e mitigação de riscos;

\textbf{VI} -- comunicar à Diretoria Executiva, ao Conselho de Administração e, se necessário, à Assembleia Geral, irregularidades ou riscos relevantes identificados;

\textbf{VII} -- convocar Assembleia Geral Extraordinária quando houver motivo grave ou omissão dos demais órgãos competentes;

\textbf{VIII} -- acompanhar o trabalho de auditoria independente, quando contratada.

\textbf{Art. 45.} O Conselho Fiscal reunir-se-á ordinariamente ao menos 2 (duas) vezes por ano e extraordinariamente sempre que necessário.

\textbf{Parágrafo único.} As reuniões instalar-se-ão com a presença da maioria de seus membros e deliberarão por maioria simples dos presentes.

\section*{TÍTULO IX -- DO PATRIMÔNIO, RECEITAS E FONTES DE RECURSOS}

\textbf{Art. 46.} O patrimônio do MoveEduca será constituído por bens móveis, imóveis, veículos, equipamentos, direitos, créditos, títulos, valores, marcas, acervos, recursos financeiros, fundos, doações, legados, subvenções e demais bens e direitos que venha a adquirir ou receber.

\textbf{Art. 47.} Constituem fontes de recursos da associação:

\textbf{I} -- contribuições de associados, mantenedores e apoiadores;

\textbf{II} -- doações, legados, auxílios, subvenções, patrocínios, apoios e repasses de pessoas físicas ou jurídicas, públicas ou privadas, nacionais ou estrangeiras;

\textbf{III} -- recursos provenientes de termos de colaboração, termos de fomento, acordos de cooperação, convênios, contratos, termos de parceria, ajustes e instrumentos congêneres;

\textbf{IV} -- receitas de editais, prêmios, chamamentos públicos, emendas, fundos, programas, projetos e incentivos legalmente admitidos;

\textbf{V} -- receitas decorrentes da prestação de serviços, consultorias, assessorias, cursos, eventos, publicações, licenciamento, desenvolvimento tecnológico, produtos institucionais e demais atividades compatíveis com as finalidades estatutárias;

\textbf{VI} -- rendimentos de aplicações financeiras, aluguéis, cessões, direitos autorais, receitas patrimoniais e outras rendas lícitas;

\textbf{VII} -- receitas decorrentes de campanhas, financiamento coletivo, bazares, comercialização de bens, materiais, produtos promocionais ou educacionais e outras iniciativas compatíveis com os objetivos sociais.

\textbf{Art. 48.} Toda receita, renda, superávit ou resultado será aplicado integralmente nas finalidades estatutárias, vedada a distribuição de lucros, excedentes ou vantagens indevidas.

\textbf{Art. 49.} A associação poderá receber recursos públicos e privados, inclusive de origem nacional ou internacional, desde que compatíveis com suas finalidades, com a legislação aplicável e com suas políticas de integridade.

\section*{TÍTULO X -- DA PRESTAÇÃO DE CONTAS, TRANSPARÊNCIA E MROSC}

\textbf{Art. 50.} O MoveEduca atuará como Organização da Sociedade Civil, nos termos da Lei nº 13.019/2014, quando preencher os requisitos legais e celebrar parcerias com a Administração Pública.

\textbf{Art. 51.} A prestação de contas da associação observará, no mínimo:

\textbf{I} -- escrituração contábil regular, de acordo com os princípios contábeis, as Normas Brasileiras de Contabilidade aplicáveis às entidades sem fins lucrativos e a legislação vigente;

\textbf{II} -- elaboração anual de relatório de atividades;

\textbf{III} -- elaboração anual de demonstrações financeiras;

\textbf{IV} -- parecer do Conselho Fiscal sobre contas, demonstrações financeiras e relatório de atividades;

\textbf{V} -- aprovação das contas pela Assembleia Geral;

\textbf{VI} -- publicidade e transparência das informações institucionais, observadas as normas de proteção de dados pessoais, sigilo legal, segurança institucional e requisitos de instrumentos de parceria;

\textbf{VII} -- guarda organizada de documentos contábeis, fiscais, trabalhistas, previdenciários, bancários, contratuais, societários, projetos, planos de trabalho e prestações de contas pelo prazo legal aplicável.

\textbf{Art. 52.} A prestação de contas de recursos públicos observará o art. 70, parágrafo único, da Constituição Federal, a Lei nº 13.019/2014, seus regulamentos, o instrumento firmado, o plano de trabalho aprovado, as normas do órgão concedente e demais legislação aplicável.

\textbf{Art. 53.} Nas parcerias com a Administração Pública, o MoveEduca observará, conforme o caso:

\textbf{I} -- chamamento público, hipóteses legais de dispensa ou inexigibilidade e demais procedimentos prévios exigidos;

\textbf{II} -- elaboração e execução de plano de trabalho;

\textbf{III} -- transparência na aplicação dos recursos;

\textbf{IV} -- segregação de contas, registros e documentos por parceria ou projeto, quando exigido;

\textbf{V} -- monitoramento, avaliação, controle de metas, indicadores, resultados e prestação de contas;

\textbf{VI} -- vedação de despesas incompatíveis com a legislação, o edital, o plano de trabalho ou o instrumento firmado;

\textbf{VII} -- restituição de valores, quando legalmente devida.

\textbf{Art. 54.} O MoveEduca poderá contratar auditoria independente quando exigida por lei, instrumento de parceria, órgão financiador, deliberação da Assembleia Geral ou recomendação fundamentada do Conselho Fiscal.

\section*{TÍTULO XI -- DA REMUNERAÇÃO, CONTRATAÇÕES E CONFLITO DE INTERESSES}

\textbf{Art. 55.} O MoveEduca poderá remunerar dirigentes que efetivamente atuem na gestão executiva ou prestem serviços específicos à associação, bem como contratar empregados, consultores, profissionais, empresas e prestadores de serviços necessários ao cumprimento de suas finalidades.

\textbf{§ 1º} A remuneração deverá ser compatível com os valores praticados pelo mercado na região de atuação, com a complexidade da função, a disponibilidade financeira da associação e a legislação aplicável.

\textbf{§ 2º} A remuneração de dirigente estatutário dependerá de aprovação do Conselho de Administração ou da Assembleia Geral, conforme definido em Regimento Interno ou política de remuneração, vedado o voto do interessado.

\textbf{§ 3º} Conselheiros fiscais e conselheiros de administração não serão remunerados pelo simples exercício da função colegiada, admitido apenas o ressarcimento de despesas comprovadas e previamente autorizadas, ou remuneração por serviço específico distinto da função de conselheiro, desde que observadas transparência, compatibilidade de mercado, aprovação do órgão competente e impedimento de voto do interessado.

\textbf{§ 4º} Quando houver recursos públicos, a remuneração e as contratações observarão o edital, o plano de trabalho, o instrumento de parceria, a Lei nº 13.019/2014, seus regulamentos e demais normas aplicáveis.

\textbf{Art. 56.} É vedada a concessão de vantagem indevida, benefício pessoal injustificado, distribuição disfarçada de resultado ou utilização de bens, recursos, oportunidades ou informações da associação em favor de associado, dirigente, conselheiro, empregado, doador, parceiro, contratado, familiar ou terceiro.

\textbf{Art. 57.} Dirigentes, conselheiros, associados, empregados, voluntários e prestadores de serviços deverão declarar conflito de interesses real, potencial ou aparente sempre que houver interesse pessoal, familiar, profissional, econômico ou institucional capaz de influenciar decisão da associação.

\textbf{§ 1º} A pessoa em conflito de interesses não poderá votar nem influenciar deliberação relacionada à matéria, devendo sua abstenção constar em ata ou registro próprio.

\textbf{§ 2º} Contratações com partes relacionadas somente poderão ocorrer quando houver justificativa técnica ou econômica, compatibilidade com valores de mercado, transparência, aprovação do órgão competente e observância de política interna, se existente.

\textbf{§ 3º} Para fins deste Estatuto, consideram-se partes relacionadas os associados, dirigentes, conselheiros, empregados em posição de decisão, seus cônjuges, companheiros, parentes até o terceiro grau, empresas ou entidades nas quais tenham participação relevante, influência ou poder de decisão.

\section*{TÍTULO XII -- DA ALTERAÇÃO ESTATUTÁRIA E DISSOLUÇÃO}

\textbf{Art. 58.} Este Estatuto poderá ser alterado por Assembleia Geral Extraordinária especialmente convocada para essa finalidade, observado o quórum previsto neste Estatuto.

\textbf{Parágrafo único.} As alterações estatutárias somente produzirão efeitos perante terceiros após registro no Cartório de Registro Civil de Pessoas Jurídicas competente, sem prejuízo da eficácia interna a partir da aprovação, quando juridicamente admissível.

\textbf{Art. 59.} O MoveEduca poderá ser dissolvido por deliberação da Assembleia Geral Extraordinária especialmente convocada para essa finalidade, quando se tornar impossível, ilícita ou inviável a continuidade de suas atividades, ou nas demais hipóteses legais.

\textbf{§ 1º} A Assembleia que deliberar a dissolução nomeará liquidante e definirá o procedimento de liquidação, respeitadas as obrigações legais, trabalhistas, fiscais, contratuais, patrimoniais e de prestação de contas.

\textbf{§ 2º} Após a liquidação, o patrimônio líquido remanescente será destinado a pessoa jurídica de direito privado sem fins lucrativos, de igual natureza, preferencialmente qualificada como Organização da Sociedade Civil e com finalidade semelhante à do MoveEduca, escolhida pela Assembleia Geral, observada a legislação aplicável.

\textbf{§ 3º} Se o MoveEduca possuir qualificação como OSCIP, ou vier a obtê-la, serão observadas adicionalmente as regras da Lei nº 9.790/1999 quanto à destinação patrimonial aplicável à respectiva qualificação.

\textbf{§ 4º} Recursos, bens ou direitos de origem pública terão destinação conforme a legislação, o instrumento jurídico respectivo e a orientação do órgão competente.

\section*{TÍTULO XIII -- DAS DISPOSIÇÕES GERAIS E TRANSITÓRIAS}

\textbf{Art. 60.} O MoveEduca poderá aprovar Regimento Interno e políticas institucionais para disciplinar seu funcionamento, governança, integridade, voluntariado, proteção de dados, compras e contratações, gestão financeira, prestação de contas, projetos, associados e demais matérias internas.

\textbf{Art. 61.} Os casos omissos serão resolvidos pela Diretoria Executiva, ad referendum da Assembleia Geral quando envolverem matéria de competência assemblear, observados a legislação aplicável, este Estatuto, o Regimento Interno e os princípios gerais de governança.

\textbf{Art. 62.} O exercício social coincidirá com o ano civil, iniciando-se em 1º de janeiro e encerrando-se em 31 de dezembro de cada ano.

\textbf{Art. 63.} Este Estatuto foi aprovado pela Assembleia Geral realizada em \textbf{[DATA DA ASSEMBLEIA]} e entrará em vigor na data de seu registro no Cartório de Registro Civil de Pessoas Jurídicas competente.

\vspace{2cm}

\begin{center}
\textbf{[CIDADE/UF]}, \textbf{[DATA DA ASSEMBLEIA]}.

\vspace{2cm}

\rule{8cm}{0.4pt}\\
\textbf{[NOME COMPLETO]}\\
Presidente do MoveEduca

\vspace{1.5cm}

\rule{8cm}{0.4pt}\\
\textbf{[NOME DO ADVOGADO]}\\
\textbf{[OAB/UF Nº]}
\end{center}

\newpage

\section*{ANEXO I -- ROL DE MEMBROS FUNDADORES, DIRETORIA E CONSELHOS}

\textbf{1. Diretoria Executiva Eleita}

\textbf{Membro(a) Fundador(a) e Presidente -- [NOME COMPLETO]}, [NACIONALIDADE], [ESTADO CIVIL], [PROFISSÃO], portador(a) do RG nº [RG], inscrito(a) no CPF sob nº [CPF], residente e domiciliado(a) em [ENDEREÇO COMPLETO], e-mail [E-MAIL], telefone [TELEFONE]. Assinatura: \rule{6cm}{0.4pt}

\vspace{1em}

\textbf{Membro(a) Fundador(a) e Vice-Presidente -- [NOME COMPLETO]}, [NACIONALIDADE], [ESTADO CIVIL], [PROFISSÃO], portador(a) do RG nº [RG], inscrito(a) no CPF sob nº [CPF], residente e domiciliado(a) em [ENDEREÇO COMPLETO], e-mail [E-MAIL], telefone [TELEFONE]. Assinatura: \rule{6cm}{0.4pt}

\vspace{1em}

\textbf{Membro(a) Fundador(a) e Diretor(a) Financeiro(a) -- [NOME COMPLETO]}, [NACIONALIDADE], [ESTADO CIVIL], [PROFISSÃO], portador(a) do RG nº [RG], inscrito(a) no CPF sob nº [CPF], residente e domiciliado(a) em [ENDEREÇO COMPLETO], e-mail [E-MAIL], telefone [TELEFONE]. Assinatura: \rule{6cm}{0.4pt}

\vspace{1em}

\textbf{Membro(a) Fundador(a) e Secretário(a)-Geral -- [NOME COMPLETO]}, [NACIONALIDADE], [ESTADO CIVIL], [PROFISSÃO], portador(a) do RG nº [RG], inscrito(a) no CPF sob nº [CPF], residente e domiciliado(a) em [ENDEREÇO COMPLETO], e-mail [E-MAIL], telefone [TELEFONE]. Assinatura: \rule{6cm}{0.4pt}

\textbf{2. Conselho de Administração}

\textbf{Membro(a) Fundador(a) e Conselheiro(a) de Administração -- [NOME COMPLETO]}, [NACIONALIDADE], [ESTADO CIVIL], [PROFISSÃO], portador(a) do RG nº [RG], inscrito(a) no CPF sob nº [CPF], residente e domiciliado(a) em [ENDEREÇO COMPLETO], e-mail [E-MAIL], telefone [TELEFONE]. Assinatura: \rule{6cm}{0.4pt}

\vspace{1em}

\textbf{Membro(a) Fundador(a) e Conselheiro(a) de Administração -- [NOME COMPLETO]}, [NACIONALIDADE], [ESTADO CIVIL], [PROFISSÃO], portador(a) do RG nº [RG], inscrito(a) no CPF sob nº [CPF], residente e domiciliado(a) em [ENDEREÇO COMPLETO], e-mail [E-MAIL], telefone [TELEFONE]. Assinatura: \rule{6cm}{0.4pt}

\vspace{1em}

\textbf{Membro(a) Fundador(a) e Conselheiro(a) de Administração -- [NOME COMPLETO]}, [NACIONALIDADE], [ESTADO CIVIL], [PROFISSÃO], portador(a) do RG nº [RG], inscrito(a) no CPF sob nº [CPF], residente e domiciliado(a) em [ENDEREÇO COMPLETO], e-mail [E-MAIL], telefone [TELEFONE]. Assinatura: \rule{6cm}{0.4pt}

\textbf{3. Conselho Fiscal}

\textbf{Membro(a) Fundador(a) e Conselheiro(a) Fiscal -- [NOME COMPLETO]}, [NACIONALIDADE], [ESTADO CIVIL], [PROFISSÃO], portador(a) do RG nº [RG], inscrito(a) no CPF sob nº [CPF], residente e domiciliado(a) em [ENDEREÇO COMPLETO], e-mail [E-MAIL], telefone [TELEFONE]. Assinatura: \rule{6cm}{0.4pt}

\vspace{1em}

\textbf{Membro(a) Fundador(a) e Conselheiro(a) Fiscal -- [NOME COMPLETO]}, [NACIONALIDADE], [ESTADO CIVIL], [PROFISSÃO], portador(a) do RG nº [RG], inscrito(a) no CPF sob nº [CPF], residente e domiciliado(a) em [ENDEREÇO COMPLETO], e-mail [E-MAIL], telefone [TELEFONE]. Assinatura: \rule{6cm}{0.4pt}

\vspace{1em}

\textbf{Membro(a) Fundador(a) e Conselheiro(a) Fiscal -- [NOME COMPLETO]}, [NACIONALIDADE], [ESTADO CIVIL], [PROFISSÃO], portador(a) do RG nº [RG], inscrito(a) no CPF sob nº [CPF], residente e domiciliado(a) em [ENDEREÇO COMPLETO], e-mail [E-MAIL], telefone [TELEFONE]. Assinatura: \rule{6cm}{0.4pt}

\newpage

\section*{CHECKLIST FINAL PARA CARTÓRIO E CNPJ}

\textbf{Para registro no Cartório de Registro Civil de Pessoas Jurídicas:}

\textbf{1.} Ata de constituição ou alteração aprovada em Assembleia, com data, local, lista de presença, mesa, ordem do dia, deliberações, aprovação do estatuto e eleição da primeira Diretoria e Conselhos.

\textbf{2.} Estatuto Social assinado pelo Presidente eleito e com visto de advogado regularmente inscrito na OAB, conforme art. 1º, § 2º, da Lei nº 8.906/1994.

\textbf{3.} Qualificação completa dos membros fundadores, dirigentes e conselheiros: nome completo, nacionalidade, estado civil, profissão, RG, CPF, endereço completo, e-mail, telefone e assinatura, quando exigida.

\textbf{4.} Lista de presença da Assembleia, com assinaturas e qualificações exigidas pelo cartório competente.

\textbf{5.} Requerimento de registro dirigido ao Cartório de Registro Civil de Pessoas Jurídicas competente.

\textbf{6.} Conferência prévia das exigências locais do cartório quanto a rubricas, reconhecimento de firma, número de vias, assinatura eletrônica, procurações e documentos pessoais.

\textbf{Para preparação do CNPJ:}

\textbf{1.} Estatuto e ata devidamente registrados no cartório competente.

\textbf{2.} DBE ou Protocolo de Transmissão da Receita Federal preenchido com a denominação, endereço, natureza jurídica, atividades econômicas e responsável perante o CNPJ.

\textbf{3.} Definição do responsável legal perante a Receita Federal, normalmente o Presidente.

\textbf{4.} Definição de CNAEs compatíveis com as finalidades educacionais, sociais, culturais, tecnológicas, de apoio institucional e demais atividades efetivamente exercidas.

\textbf{5.} Documentos pessoais e comprovante de endereço do responsável, se exigidos.

\textbf{6.} Verificação de inscrições municipais, alvarás, licenças, certificado digital, conta bancária, cadastro em plataformas públicas e obrigações acessórias aplicáveis à atuação concreta da associação.

\textbf{Atenção:} antes do protocolo, recomenda-se revisão final por advogado inscrito na OAB e por profissional contábil, especialmente para adequação ao cartório competente, CNAEs, endereço, qualificação dos dirigentes, exigências locais e eventual estratégia de captação de recursos públicos ou privados.

\end{document}
